\documentclass[11pt]{article}
\usepackage[margin=1in]{geometry}
\usepackage{hyperref}
\usepackage{enumitem}
\title{Digital Coloring and Tool Routines}
\author{Piter Garcia}
\date{\today}
\begin{document}
\maketitle

\section*{Grade and Class Match}
\textbf{Grade:} Kindergarten\\
\textbf{Likely blocks:} Day 1 · 12:25-1:15 · Justinger; Day 2 · 12:25-1:15 · Kesys; Day 4 · 12:25-1:15 · Pum

\section*{Lesson Focus}
Device routines, calm transitions, and simple digital creativity.

\section*{Learning Targets}
\begin{itemize}[leftmargin=*]
\item Students open and use a digital coloring tool with minimal support.
\item Students follow device-handling and transition routines safely.
\item Students describe one tool or action they used during the task.
\end{itemize}

\section*{Materials}
\begin{itemize}[leftmargin=*]
\item Student devices
\item Digital coloring app or site
\item Projected model example
\end{itemize}

\section*{Suggested Lesson Flow}
\begin{enumerate}[leftmargin=*]
\item Open with the exact device, tool, or routine students need in order to start successfully.
\item Model one example task before students begin independent or partner work.
\item Give students time to build, code, design, or document their work while circulating for support.
\item Close with a short share-out, reflection, or debugging discussion.
\end{enumerate}

\section*{Source Notes}
This lesson packet is enriched with transcript-derived lesson notes, the school schedule roster, and official starting-point resources. See the paired Markdown guide for clickable links.

\bibliographystyle{plain}
\bibliography{digital-coloring-and-tool-routines}
\end{document}
